\documentclass[11pt,a4paper]{article}
\newcommand{\resultsdir}{../../results/task2}
% Common preamble of the task reports (compile with XeLaTeX).
\usepackage{fontspec}
\usepackage{polyglossia}
\setmainlanguage{polish}
\usepackage{amsmath,amssymb}
\usepackage[a4paper,margin=2.3cm]{geometry}
\usepackage{graphicx}
\usepackage{booktabs}
\usepackage{longtable}
\usepackage{array}
\usepackage{multirow}
\usepackage{adjustbox}
\usepackage{makecell}
\usepackage{placeins}
\usepackage{float}
\usepackage[font=small,labelfont=bf]{caption}
\usepackage{subcaption}
\usepackage{xcolor}
\usepackage{enumitem}
\usepackage[hidelinks]{hyperref}

% Without the texlive-lang-polish package XeLaTeX would hyphenate Polish
% words with English patterns, so hyphenation is switched off and looser
% inter-word spacing is allowed instead. With Polish patterns installed the
% next two penalties can be removed.
\hyphenpenalty=10000
\exhyphenpenalty=10000
\tolerance=3000
\emergencystretch=3em
\setlist{itemsep=2pt,topsep=4pt}
\setlength{\parskip}{3pt}
\renewcommand{\arraystretch}{1.12}

\newcommand{\tab}[1]{\input{\resultsdir/tables/#1}}
% table scaled down to the text width only when it is wider
\newcommand{\fittab}[1]{\adjustbox{max width=\textwidth}{\tab{#1}}}
\newcommand{\fig}[1]{\resultsdir/figures/#1}
\newcommand{\code}[1]{\texttt{#1}}
% left-aligned paragraph column (no stretched spaces in narrow cells)
\newcolumntype{L}[1]{>{\raggedright\arraybackslash}p{#1}}

% \maketitlepage{task title}{subtitle}
\newcommand{\maketitlepage}[2]{%
  \begin{titlepage}
    \centering
    {\large Politechnika Łódzka\\Wydział Fizyki Technicznej, Informatyki i Matematyki Stosowanej\par}
    \vspace{3.2cm}
    {\large Analiza danych złożonych z detekcją wyjątków\par}
    \vspace{0.8cm}
    {\LARGE\bfseries #1\par}
    \vspace{0.4cm}
    {\large #2\par}
    \vfill
    {\large Jędrzej Kostyk\par}
    {nr indeksu 264339\par}
    \vspace{1.2cm}
    {Rok akademicki 2026/2027\par}
  \end{titlepage}
}


\begin{document}
\maketitlepage{Zadanie 2}{Detekcja wyjątków w danych wielowymiarowych:\\wpływ reprezentacji danych}
\tableofcontents
\newpage

\section{Cel i zakres}

Celem zadania jest sprawdzenie, czy detektor wyjątków działa lepiej, gdy
zamiast surowych cech dostanie inną \emph{reprezentację} danych, czyli te
same obserwacje opisane innym zestawem liczb. Porównano cztery
reprezentacje: dane oryginalne, PCA, UMAP i przestrzeń ukrytą autoenkodera.
W każdej z nich uruchomiono te same trzy detektory z zadania~1 -- Isolation
Forest, LOF i One-Class SVM. PCA i autoenkoder dają też własny wynik
anomalii (błąd rekonstrukcji), który oceniono jako czwarty detektor.

Eksperymenty przeprowadzono na czterech zbiorach danych o 21--274 cechach.
Dla każdej reprezentacji zbadano jej parametry (np. liczbę składowych PCA),
co daje do 33 wariantów reprezentacji i 15 ustawień detektorów na każdy
podział danych. Cały kod znajduje się w katalogach \code{common/}
i \code{task2\_representations/} repozytorium.

\FloatBarrier
\section{Dane}

W tej sekcji przedstawiono cztery zbiory danych użyte w zadaniu 2 (te same
zbiory wykorzystuje zadanie 3).

\FloatBarrier
\subsection{Zbiory danych}

Zadanie wymaga danych z wieloma cechami, dlatego wybrano cztery gotowe zbiory
z benchmarków ADBench i ODDS (tabela~\ref{tab:datasets}). Etykiety anomalii
pochodzą od autorów benchmarków; zbiory nie były w żaden sposób zmieniane.
Cardio występuje też w zadaniu~1, co pozwala porównać oba zadania.

\begin{table}[H]
  \centering\small
  \caption{Zbiory danych użyte w zadaniu.}
  \label{tab:datasets}
  \fittab{datasets}
\end{table}

Zbiory różnią się tym, co utrudnia detekcję:
\begin{itemize}
  \item \textbf{Arrhythmia} ma więcej cech (274) niż obserwacji uczących
  (226), a 17 cech jest stałych;
  \item \textbf{WDBC} ma tylko 10 anomalii, więc w części testowej są 2--3;
  \item \textbf{Campaign} jest duży (41 tys. obserwacji), a część cech to
  zakodowane kategorie (np. zawód klienta zapisany jako kilka kolumn 0/1).
\end{itemize}

Tabela~\ref{tab:samples} pokazuje po jednej obserwacji normalnej i jednej
anomalii z każdego zbioru danych. Benchmarki udostępniają cechy bez nazw,
dlatego oznaczono je numerami.

\begin{table}[H]
  \centering\small
  \setlength{\tabcolsep}{5pt}
  \caption{Przykładowe obserwacje (pierwsze pięć cech; ,,\dots'' -- pozostałe cechy).}
  \label{tab:samples}
  \textbf{Cardio}\\[2pt]\fittab{sample_cardio}\\[6pt]
  \textbf{Arrhythmia}\\[2pt]\fittab{sample_arrhythmia}\\[6pt]
  \textbf{WDBC}\\[2pt]\fittab{sample_wdbc}\\[6pt]
  \textbf{Campaign}\\[2pt]\fittab{sample_campaign}
\end{table}

\FloatBarrier
\section{Protokół eksperymentalny}
\label{sec:protocol}

Ta sekcja opisuje, jak przygotowano eksperyment, żeby wyniki były uczciwe
i porównywalne z zadaniem~1. Protokół jest taki sam jak w zadaniu~1; nowa
jest tylko reprezentacja, która też uczy się wyłącznie na części uczącej.

\paragraph{Podział na trzy części.} Każdy zbiór danych dzielimy losowo na
część uczącą (50\,\%), walidacyjną (25\,\%) i testową (25\,\%), z takim
samym udziałem anomalii w każdej części. Przykład: Arrhythmia ma 452
obserwacje, więc 226 trafia do uczenia, 113 do walidacji i 113 do testu.
\begin{itemize}
  \item \textbf{uczenie} -- na nim powstaje reprezentacja (np. osie PCA)
  i uczy się detektor, \emph{bez} informacji, które obserwacje są anomaliami;
  \item \textbf{walidacja} -- na niej wybieramy parametry reprezentacji,
  ustawienia detektora i próg decyzji;
  \item \textbf{test} -- oceniamy na nim wynik tylko raz, na końcu.
\end{itemize}

\paragraph{Skalowanie i reprezentacja tylko z części uczącej.} Każdą cechę
standaryzujemy (odejmujemy średnią i dzielimy przez odchylenie standardowe),
licząc średnią i odchylenie tylko na części uczącej. Tak samo reprezentacja:
np. osie PCA wyznaczamy z 226 obserwacji uczących Arrhythmii, a potem tymi
samymi osiami przekształcamy walidację i test. Na koniec każdą reprezentację
mnożymy przez jedną liczbę, tak by jej współrzędne miały średnio taki sam
rozrzut jak w danych oryginalnych. Nie zmienia to kształtu danych, ale
sprawia, że parametr detektora (np. szerokość jądra OCSVM) znaczy to samo
w każdej reprezentacji.

\paragraph{Ziarno losowe i wiele podziałów.} \emph{Ziarno} to liczba
startowa generatora losowego: to samo ziarno daje zawsze ten sam podział.
Jeden podział to loteria -- WDBC ma w teście 2--3 anomalie, więc jedno
trafienie zmienia wynik o kilkadziesiąt punktów procentowych. Dlatego Cardio
i Arrhythmię dzielimy 10 razy, WDBC 30 razy, a Campaign 5 razy (koszt
obliczeń, niżej). Raportujemy średnią i odchylenie standardowe po
podziałach. Te same ziarna dają te same podziały w zadaniach 1--3.

\paragraph{Wybór ustawień i progu.} Detektor daje każdej obserwacji
\emph{wynik anomalii} -- liczbę tym większą, im bardziej obserwacja jest
podejrzana. Parametry reprezentacji i detektora wybieramy razem: to
połączenie, które ma najwyższą średnią PR-AUC na walidacji. Następnie na
walidacji wybieramy próg wyniku, który daje najwyższe F1.

\paragraph{Miary jakości.} Te same co w zadaniu~1:
\begin{itemize}
  \item \textbf{Precision} -- ile zgłoszonych obserwacji to naprawdę anomalie;
  \item \textbf{Recall} -- ile prawdziwych anomalii wykryto;
  \item \textbf{F1} -- jedna liczba łącząca Precision i Recall, wysoka tylko
  wtedy, gdy obie są wysokie;
  \item \textbf{ROC-AUC} -- szansa, że losowa anomalia dostanie wyższy wynik
  niż losowa obserwacja normalna (0,5 = zgadywanie);
  \item \textbf{PR-AUC} -- jak wysoko na liście najbardziej podejrzanych
  obserwacji są anomalie (od 0 do 1; zgadywanie daje wynik równy udziałowi
  anomalii, np. 0,10 dla Cardio). To główna miara porównań.
\end{itemize}
Mierzono też czas: przygotowania reprezentacji oraz uczenia i oceny detektora.

\paragraph{Ograniczenie obliczeń dla Campaign.} Campaign ma 20\,594
obserwacje uczące, a jeden model OCSVM uczy się na nich ok. 40~s, UMAP
25--40~s, a autoenkoder kilkadziesiąt sekund. Dlatego dla tego zbioru,
\textbf{bez zmniejszania danych}, użyto 5 zamiast 10 podziałów, 6 zamiast 18
wariantów UMAP i OCSVM bez $\nu=0{,}5$. Miary struktury reprezentacji
(sekcja~\ref{sec:methods}) liczono na stałym losowym podzbiorze 2000 punktów
testowych; detekcję oceniano zawsze na całej części testowej.

\FloatBarrier
\section{Reprezentacje i detektory}
\label{sec:methods}

Tabela~\ref{tab:reps} opisuje, jak działa każda reprezentacja i jakie
wartości jej parametrów sprawdzono, a tabela~\ref{tab:detectors} --
detektory uruchamiane w każdej z nich.

\begin{table}[H]
  \centering\small
  \caption{Reprezentacje danych i sprawdzone wartości parametrów.}
  \label{tab:reps}
  \begin{tabular}{L{2.4cm}L{7.6cm}L{4.9cm}}
    \toprule
    Reprezentacja & Jak działa & Parametry (sprawdzone wartości) \\
    \midrule
    dane oryginalne & Wszystkie cechy po standaryzacji; punkt odniesienia. &
    -- \\
    \addlinespace
    PCA (analiza składowych głównych) & Zastępuje cechy ich kombinacjami
    (\emph{składowymi}) uporządkowanymi od tej, wzdłuż której dane są
    najbardziej rozrzucone. Zostawiamy $k$ pierwszych składowych, a resztę
    odrzucamy. \emph{Zachowana wariancja} mówi, jaką część rozrzutu danych
    zachowuje $k$ składowych (1 = całą). &
    liczba składowych $k$: 1, 2, 3, 5, 8, 10, 15, 20 \\
    \addlinespace
    UMAP & Nieliniowo ,,spłaszcza'' dane do $m$ wymiarów tak, by punkty
    bliskie w oryginale pozostały blisko siebie. Nowe punkty umieszcza obok ich
    najbliższych sąsiadów z części uczącej. &
    wymiar $m$: 2, 5, 10; liczba sąsiadów nn (lokalna albo szersza struktura):
    5, 15, 50; minimalny odstęp punktów md: 0; 0,5 \\
    \addlinespace
    autoenkoder & Sieć neuronowa, która ściska dane do $z$ liczb
    (\emph{przestrzeń ukryta}) i próbuje odtworzyć z nich wszystkie cechy.
    Jako reprezentację bierzemy te $z$ liczb. &
    wymiar przestrzeni ukrytej $z$: 2, 4, 8, 16; dla $z=8$ także
    dwie inne architektury \\
    \addlinespace
    błąd rekonstrukcji (PCA, autoenkoder) & Nie reprezentacja, lecz
    dodatkowy detektor: jak bardzo obserwacja odtworzona z $k$ składowych
    lub z $z$ liczb różni się od oryginału. Obserwacja niepasująca do
    typowej struktury danych odtwarza się źle. &
    jak dla PCA i autoenkodera \\
    \bottomrule
  \end{tabular}
\end{table}

\begin{table}[H]
  \centering\small
  \caption{Detektory uruchamiane w każdej reprezentacji ($d$ -- liczba
  wymiarów reprezentacji).}
  \label{tab:detectors}
  \begin{tabular}{L{2.4cm}L{7.6cm}L{4.9cm}}
    \toprule
    Detektor & Jak działa & Parametry (sprawdzone wartości) \\
    \midrule
    Isolation Forest & Anomalię da się odciąć od reszty po niewielu losowych
    cięciach. & liczba drzew $T=200$ (w zadaniu~1 nie miała wpływu) \\
    \addlinespace
    LOF & Punkt wyraźnie rzadszy niż jego sąsiedzi jest podejrzany. &
    liczba sąsiadów $k_{LOF}$: 5, 10, 20, 35, 50 \\
    \addlinespace
    One-Class SVM & Granica wokół obszaru danych; punkty poza nią są
    anomaliami. Jądro RBF. & $\nu$ (limit odsetka punktów uczących poza
    granicą): 0,05; 0,1; 0,5; $\gamma$ (szerokość jądra): 0,1/$d$, 1/$d$,
    10/$d$ \\
    \bottomrule
  \end{tabular}
\end{table}

Ustawienia detektora wybieramy osobno w każdej reprezentacji, bo najlepsze
$k_{LOF}$ czy $\gamma$ zależą od kształtu danych. Wspólne ustawienia
faworyzowałyby reprezentację, do której je dobrano.

\paragraph{Architektura autoenkodera.} Sieć ma kształt lejka: $d \to 64
\to 32 \to z \to 32 \to 64 \to d$, gdzie $d$ to liczba cech. Uzasadnienie:
\begin{itemize}
  \item środek węższy niż wejście ($z<d$) zmusza sieć do kompresji, więc nie
  może po prostu przepisać danych, tylko uczy się ich typowej struktury;
  \item dwie warstwy ukryte pozwalają uchwycić zależności nieliniowe, których
  PCA nie widzi;
  \item małe zbiory (183--915 obserwacji uczących poza Campaign) łatwo się
  ,,uczą na pamięć'', dlatego uczenie zatrzymujemy, gdy błąd na 15\,\% części
  uczącej odłożonych z boku przestaje maleć (maks. 400 epok, czyli przejść
  przez dane).
\end{itemize}
Sieć uczy się na danych z anomaliami, tak jak w praktyce, gdy etykiet nie ma.
Dwie inne architektury ($d\to32\to8$ i $d\to128\to64\to8$) porównano
w sekcji~\ref{sec:ae}.

\paragraph{Miary struktury reprezentacji.} Dla każdej reprezentacji
policzono na części testowej dwie miary, które nie zależą od detektora:
\begin{itemize}
  \item \textbf{zaufanie sąsiedztw} (\emph{trustworthiness}, od 0 do 1) --
  czy sąsiedzi punktu w nowej reprezentacji byli też jego sąsiadami
  w danych oryginalnych (1 = brak ,,fałszywych'' sąsiadów);
  \item \textbf{separacja klas} (\emph{silhouette}, od $-1$ do 1) -- czy
  anomalie leżą bliżej siebie niż obserwacji normalnych; 0 = klasy wymieszane.
\end{itemize}
Etykiety służą tu tylko do analizy, nie do budowy reprezentacji. Recall
liczono też osobno dla anomalii \emph{globalnych} (leżących dalej od środka
danych normalnych niż 99\,\% obserwacji normalnych) i \emph{lokalnych}
(pozostałych), tak jak w zadaniu~1.

\FloatBarrier
\section{Wyniki}

W tej sekcji najpierw zestawiono wszystkie pary reprezentacja--detektor,
potem omówiono każdą reprezentację osobno w tym samym układzie, a na końcu
porównano je ze sobą i z zadaniem~1.

\FloatBarrier
\subsection{Przegląd: wszystkie reprezentacje i detektory}

Rysunek~\ref{fig:matrix} pokazuje jakość każdej pary reprezentacja--detektor
na każdym zbiorze danych; pełne wyniki są w aneksie~\ref{app:full}.

\begin{figure}[H]
  \centering
  \includegraphics[width=0.95\textwidth]{\fig{pr_auc_matrix}}
  \caption{PR-AUC każdej pary reprezentacja--detektor (ustawienia wybrane na
  walidacji); ciemniejsze pole -- lepszy wynik, ,,--'' -- nie dotyczy.}
  \label{fig:matrix}
\end{figure}

\textbf{Zmiana reprezentacji może zmienić wynik detektora o ponad 0,2
PR-AUC, w obie strony}:
\begin{itemize}
  \item \textbf{WDBC: PCA $\Rightarrow$ LOF 0,71 $\rightarrow$ 0,92, OCSVM
  0,70 $\rightarrow$ 0,93};
  \item \textbf{Cardio: UMAP $\Rightarrow$ OCSVM 0,60 $\rightarrow$ 0,84}
  (najlepszy wynik na tym zbiorze), ale LOF 0,31 $\rightarrow$ 0,24;
  \item \textbf{Arrhythmia: UMAP $\Rightarrow$ IForest 0,58 $\rightarrow$
  0,33}, a autoenkoder i PCA poprawiają LOF i OCSVM o 0,09--0,14;
  \item \textbf{Campaign: żadna reprezentacja nie pomaga} -- wszystkie pary
  mają 0,14--0,30 przy zgadywaniu 0,11.
\end{itemize}

Tabela~\ref{tab:spread} pokazuje, który parametr reprezentacji ma największy
wpływ: o ile zmienia się PR-AUC, gdy zmieniamy tylko ten parametr (resztę,
łącznie z detektorem, wybiera walidacja).

\begin{table}[H]
  \centering\small
  \caption{Wpływ parametrów reprezentacji: różnica PR-AUC między najlepszą
  a najgorszą wartością parametru (,,--'' -- na Campaign badano jedną wartość).}
  \label{tab:spread}
  \fittab{param_spread}
\end{table}

\textbf{Największy wpływ mają liczba składowych PCA (do 0,21 na WDBC)
i liczba sąsiadów UMAP (0,31 na Cardio)}; wymiar UMAP, odstęp md i wymiar
autoenkodera zmieniają wynik najwyżej o 0,08.

Rysunki ,,co wykryto'' w kolejnych podsekcjach pokazują, co najlepszy (według
walidacji) detektor w danej reprezentacji wykrył w części testowej jednego
podziału (ziarno 0). Wszystkie mapy używają tego samego rzutu -- dwóch
pierwszych składowych PCA danych oryginalnych -- więc można je porównywać
punkt po punkcie. Kolory oznaczają: \textbf{TP} -- wykryta anomalia,
\textbf{FN} -- przeoczona anomalia, \textbf{FP} -- fałszywy alarm; szare
punkty to obserwacje normalne słusznie niezgłoszone. Część testowa WDBC
w tym podziale ma tylko 2 anomalie.

\FloatBarrier
\subsection{Dane oryginalne (punkt odniesienia)}

Tu pokazano, co wykrywają detektory na samych cechach, bez zmiany
reprezentacji.

\begin{figure}[H]
  \centering
  \includegraphics[width=\textwidth]{\fig{detected_original}}
  \caption{Co wykrył najlepszy detektor w danych oryginalnych na każdym
  zbiorze danych (w nawiasie -- detektor wybrany na walidacji).}
  \label{fig:det_original}
\end{figure}

\textbf{Na Cardio OCSVM wykrywa 36 z 44 anomalii, ale z 36 fałszywymi
alarmami; na Campaign IForest przeocza 763 z 1160 anomalii}
(rys.~\ref{fig:det_original}). Anomalie Campaign leżą w tych samych
obszarach co klienci, którzy lokaty nie założyli, więc żaden detektor nie
odróżnia ich po położeniu.

\FloatBarrier
\subsection{PCA}

Sprawdzono, jak liczba zostawionych składowych $k$ wpływa na wykrywanie.

\begin{figure}[H]
  \centering
  \includegraphics[width=\textwidth]{\fig{pca_components}}
  \caption{PCA: PR-AUC w zależności od liczby składowych $k$; linia szara --
  zachowana wariancja (ta sama skala 0--1).}
  \label{fig:pca}
\end{figure}

Zależność między zachowaną wariancją a jakością detekcji \textbf{nie jest
monotoniczna} (rys.~\ref{fig:pca}):
\begin{itemize}
  \item \textbf{WDBC: $k=1$ (33\,\% wariancji) $\Rightarrow$ OCSVM 0,93, LOF
  0,92} -- pierwsza składowa opisuje ogólną wielkość jądra komórkowego, czyli
  właśnie to, czym guzy złośliwe różnią się od łagodnych; każda kolejna
  składowa dodaje tylko zmienność wśród guzów łagodnych.
  \item \textbf{Cardio: $k=1$ $\Rightarrow$ LOF 0,58} (w danych
  oryginalnych 0,31); OCSVM jest najlepszy przy $k=1$--5.
  \item \textbf{Arrhythmia: $k=5$--10 $\Rightarrow$ LOF 0,59}, choć to
  tylko 32--47\,\% wariancji. Przy 274 cechach odległości w danych
  oryginalnych zagłusza szum setek słabych cech, a redukcja go usuwa. Przy
  $k\le2$ wynik spada (0,34--0,40), bo znika też informacja o anomaliach.
  \item \textbf{Campaign: $k$ prawie nie ma wpływu} (IForest 0,24--0,30).
\end{itemize}
Zachowana wariancja mierzy, jak wiernie odtwarzamy \emph{typowe} dane, a nie
to, czy zostaje informacja odróżniająca \emph{nietypowe}. Dlatego $k$ trzeba
wybierać na walidacji, a nie według progu wariancji.

\begin{table}[H]
  \centering\small
  \caption{PCA: LOF i błąd rekonstrukcji przy różnej liczbie składowych $k$
  (,,wszystkie'' -- dane oryginalne; Recall osobno dla anomalii globalnych
  i lokalnych).}
  \label{tab:pca_types}
  \fittab{pca_types}
\end{table}

Tabela~\ref{tab:pca_types} pokazuje, co redukcja robi z typami anomalii.
\textbf{Arrhythmia: $k=2$ $\Rightarrow$ Recall globalnych anomalii LOF 0,87
$\rightarrow$ 0,54}: odchylenie takiej anomalii rozkłada się na wiele cech,
a rzut na dwie składowe je usuwa. To samo odchylenie widać jednak w błędzie
rekonstrukcji, który wykrywa 0,86--0,99 anomalii globalnych na każdym zbiorze
danych. Ma przy tym niską trafność zgłoszeń (PR-AUC 0,24--0,47), bo równie
źle odtwarza się wiele obserwacji normalnych.

\begin{figure}[H]
  \centering
  \includegraphics[width=\textwidth]{\fig{detected_pca}}
  \caption{Co wykrył najlepszy detektor w reprezentacji PCA na każdym zbiorze
  danych.}
  \label{fig:det_pca}
\end{figure}

W porównaniu z danymi oryginalnymi \textbf{PCA przede wszystkim ogranicza
fałszywe alarmy}: na Cardio OCSVM zgłasza 18 zamiast 36, a na Campaign
IForest 592 zamiast 870 (rys.~\ref{fig:det_pca}), ale kosztem części
wykrytych anomalii (Cardio: 29 zamiast 36, Campaign: 327 zamiast 397).

\FloatBarrier
\subsection{UMAP}
\label{sec:umap}

Sprawdzono wpływ liczby sąsiadów nn, wymiaru $m$ i odstępu md na
wykrywanie.

\begin{figure}[H]
  \centering
  \includegraphics[width=\textwidth]{\fig{umap_params}}
  \caption{UMAP: PR-AUC w zależności od liczby sąsiadów nn ($m$ i md
  wybrane na walidacji); liczba sąsiadów działa w przeciwnych kierunkach na
  Cardio i WDBC.}
  \label{fig:umap}
\end{figure}

\begin{itemize}
  \item \textbf{Liczba sąsiadów nn -- Cardio: 5 $\rightarrow$ 50
  $\Rightarrow$ OCSVM 0,53 $\rightarrow$ 0,84; WDBC: odwrotnie, IForest 0,66
  $\rightarrow$ 0,12} (rys.~\ref{fig:umap}). Anomalie Cardio to duża, spójna
  grupa, którą UMAP przy szerszym kontekście odsuwa od reszty danych. W WDBC
  jest tylko 5 anomalii w części uczącej i szeroki kontekst ,,wciąga'' je
  między obserwacje normalne.
  \item \textbf{Wymiar $m$ i odstęp md} zmieniają wynik najwyżej o 0,08
  (tabela~\ref{tab:spread}); $m=2$ wystarcza do rysunku, ale do detekcji
  zwykle lepsze jest 5--10.
\end{itemize}

\textbf{Dlaczego UMAP zwykle szkodzi?} UMAP z założenia wyrównuje gęstość:
każdy punkt dostaje ,,tyle samo'' sąsiadów, bez względu na to, jak daleko
oni leżą. Odosobniona anomalia zostaje dociągnięta do najbliższej grupy,
a mała grupa anomalii staje się zwartym skupiskiem. LOF, który szuka
punktów rzadszych od sąsiadów, przestaje je widzieć (Cardio: 0,31
$\rightarrow$ 0,24). \textbf{Dlaczego pomaga OCSVM na Cardio?} OCSVM nie
potrzebuje, żeby anomalie były rzadkie -- wystarczy, że leżą poza główną masą
danych, a UMAP odsuwa grupę zapisów patologicznych na brzeg.

\begin{figure}[H]
  \centering
  \includegraphics[width=\textwidth]{\fig{detected_umap}}
  \caption{Co wykrył najlepszy detektor w reprezentacji UMAP na każdym zbiorze
  danych.}
  \label{fig:det_umap}
\end{figure}

\textbf{Na Cardio OCSVM w UMAP ma najlepszą równowagę: 34 z 44 anomalii przy
10 fałszywych alarmach} (w danych oryginalnych 36 i 36). Na Arrhythmii LOF
w UMAP wykrywa tylko 4 z 16 anomalii (rys.~\ref{fig:det_umap}).

\FloatBarrier
\subsection{Autoenkoder}
\label{sec:ae}

Sprawdzono wpływ wymiaru przestrzeni ukrytej $z$ i architektury sieci.

\begin{figure}[H]
  \centering
  \includegraphics[width=\textwidth]{\fig{ae_latent}}
  \caption{Autoenkoder: PR-AUC detektorów w przestrzeni ukrytej i błędu
  rekonstrukcji w zależności od wymiaru $z$; wymiar ma mały wpływ.}
  \label{fig:ae_latent}
\end{figure}

\textbf{Wymiar $z$ zmienia wynik najwyżej o 0,07} (rys.~\ref{fig:ae_latent},
tabela~\ref{tab:spread}). Autoenkoder wyraźnie pomaga tylko na Arrhythmii
(LOF 0,46 $\rightarrow$ 0,60), gdzie problemem jest nadmiar cech, i działa
tam podobnie jak PCA (0,59). Na pozostałych zbiorach zyskuje mniej niż PCA
(np. LOF na WDBC: 0,71 $\rightarrow$ 0,75, w PCA 0,92) albo traci: sieć uczy
się odtwarzać dane, a nie zachowywać odległości, na których opierają się
detektory.

\begin{table}[H]
  \centering\small
  \caption{Porównanie architektur autoenkodera przy $z=8$: PR-AUC błędu
  rekonstrukcji, PR-AUC najlepszego detektora w przestrzeni ukrytej (w
  nawiasie -- detektor), średnia liczba epok do zatrzymania i czas uczenia.}
  \label{tab:arch}
  \fittab{ae_architecture}
\end{table}

\textbf{Architektura zmienia wynik najwyżej o 0,05, z wyjątkiem Campaign},
gdzie płytsza sieć $d$-32-8 ma za małą pojemność (0,21 wobec 0,28;
tabela~\ref{tab:arch}). Na Campaign uczenie zawsze dochodzi do limitu 400
epok, czyli sieć nie osiągnęła jeszcze najlepszego wyniku. Na Cardio
błąd rekonstrukcji wykrywa słabo (0,25--0,26): 88 zapisów patologicznych
w części uczącej tworzy własną grupę, którą sieć uczy się odtwarzać.

\begin{figure}[H]
  \centering
  \includegraphics[width=\textwidth]{\fig{detected_ae}}
  \caption{Co wykrył najlepszy detektor w reprezentacji autoenkodera na każdym
  zbiorze danych.}
  \label{fig:det_ae}
\end{figure}

\textbf{Na Campaign najlepszy według walidacji jest błąd rekonstrukcji
autoenkodera: wykrywa 497 z 1160 anomalii (IForest w danych oryginalnych 397),
ale zgłasza też 905 fałszywych alarmów} (rys.~\ref{fig:det_ae}). Na Cardio
OCSVM w przestrzeni ukrytej wykrywa tyle samo anomalii co w danych
oryginalnych (36), z mniejszą liczbą fałszywych alarmów (27 wobec 36).

\FloatBarrier
\subsection{Co wykryto: porównanie reprezentacji}

Tabela~\ref{tab:detected} zbiera liczby z czterech map ,,co wykryto''.

\begin{table}[H]
  \centering\small
  \caption{Liczby z rysunków ,,co wykryto'' (część testowa, ziarno 0):
  TP / FN / FP, czyli wykryte anomalie / przeoczone anomalie / fałszywe
  alarmy; w nawiasie -- detektor wybrany na walidacji.}
  \label{tab:detected}
  \fittab{detected_counts}
\end{table}

\textbf{Reprezentacja zmienia głównie liczbę fałszywych alarmów, a mniej
liczbę wykrytych anomalii}: na Cardio wszystkie cztery reprezentacje wykrywają
29--36 z 44 anomalii, ale fałszywych alarmów jest od 10 (UMAP) do 36 (dane
oryginalne).

\FloatBarrier
\subsection{Wizualizacje 2D i separacja klas}

Ta część sprawdza, czy reprezentacja, w której anomalie \emph{widać}
najlepiej, daje też najlepszą detekcję.

\begin{figure}[H]
  \centering
  \includegraphics[width=\textwidth]{\fig{embeddings_2d}}
  \caption{Dwuwymiarowe reprezentacje części testowej (ziarno 0): PCA,
  UMAP i autoenkoder; ładny obraz nie oznacza dobrej detekcji.}
  \label{fig:emb}
\end{figure}

\textbf{Obraz 2D może mylić w obie strony} (rys.~\ref{fig:emb}). UMAP daje
najbardziej ,,czytelne'' obrazy: Campaign rozbija na kilkadziesiąt skupisk
(odpowiadających kombinacjom kategorii), a Arrhythmię na wyraźne grupy, ale
anomalie są w nich rozproszone i detekcja jest słaba. Na Cardio UMAP odsuwa
anomalie na brzeg i to daje najlepszy wynik. PCA na 2 składowe dobrze
oddziela anomalie WDBC, ale w Arrhythmii (17\,\% wariancji) ukrywa
większość struktury.

\begin{table}[H]
  \centering\small
  \caption{Zaufanie sąsiedztw w UMAP i separacja klas w każdej reprezentacji
  (od -- do po wariantach reprezentacji).}
  \label{tab:structure}
  \fittab{structure}
\end{table}

\begin{figure}[H]
  \centering
  \includegraphics[width=\textwidth]{\fig{separation_vs_detection}}
  \caption{Separacja klas a PR-AUC dla wszystkich wariantów reprezentacji;
  $r$ -- współczynnik korelacji (1 = pełna zgodność, 0 = brak związku).}
  \label{fig:sep}
\end{figure}

\textbf{Separacja klas zgadza się z jakością detekcji na Arrhythmii
($r=0{,}95$) i WDBC ($r=0{,}89$), ale nie na Campaign ($r=0{,}28$) i Cardio
($r=-0{,}22$)} (rys.~\ref{fig:sep}). Na Cardio najlepiej wykrywa UMAP, choć
ma niższą separację niż dane oryginalne (0,22--0,30 wobec 0,33;
tabela~\ref{tab:structure}). Separacja mierzy, czy anomalie tworzą osobną
grupę, a nie czy leżą poza obszarem danych normalnych. UMAP najsłabiej
zachowuje sąsiedztwa na Arrhythmii (0,62--0,80), gdzie 274 cechy ściska do
najwyżej 10.

\FloatBarrier
\subsection{Anomalie globalne i lokalne}

Ta część sprawdza, który typ anomalii zyskuje, a który traci na zmianie
reprezentacji.

\begin{table}[H]
  \centering\small
  \caption{Recall anomalii globalnych i lokalnych najlepszego (według
  walidacji) detektora w każdej reprezentacji.}
  \label{tab:types}
  \fittab{recall_by_type}
\end{table}

\textbf{UMAP obniża wykrywanie anomalii globalnych na każdym zbiorze danych}
(Arrhythmia 0,84 $\rightarrow$ 0,66, Campaign 0,83 $\rightarrow$ 0,63,
Cardio 0,97 $\rightarrow$ 0,85; tabela~\ref{tab:types}), bo traci informację
o tym, \emph{jak daleko} anomalia leży. Anomalie lokalne zyskują na
reprezentacji, która usuwa szum: PCA na WDBC (0,58 $\rightarrow$ 0,76)
i UMAP na Cardio (0,62 $\rightarrow$ 0,75). Żadna para reprezentacja--detektor
nie jest najlepsza dla obu typów na wszystkich zbiorach danych.

\FloatBarrier
\subsection{Jakość a czas działania}

Ta część zestawia jakość i czas na wszystkich czterech zbiorach danych; dla
każdej reprezentacji użyto detektora wybranego na walidacji.

\begin{table}[H]
  \centering\small
  \caption{Jakość i czas reprezentacji na czterech zbiorach danych. Miejsce --
  pozycja reprezentacji według PR-AUC na danym zbiorze danych (1 =
  najlepsza), uśredniona po zbiorach.}
  \label{tab:time}
  \fittab{time_quality}
\end{table}

\begin{figure}[H]
  \centering
  \includegraphics[width=\textwidth]{\fig{time_vs_quality}}
  \caption{(a) Średnia jakość a łączny czas (kreska -- od najkrótszego do
  najdłuższego czasu na czterech zbiorach danych); (b) czas przygotowania
  reprezentacji a liczba punktów uczących (obie osie logarytmiczne).}
  \label{fig:time}
\end{figure}

\textbf{PCA jest jednocześnie najlepsza i najtańsza: średnio PR-AUC 0,617
i 1,5 miejsce, przy 5~ms na przygotowanie reprezentacji}
(tabela~\ref{tab:time}). UMAP i autoenkoder kosztują tysiące razy więcej
(średnio 7,6~s i 15,8~s, na Campaign do 27~s i 49~s; rys.~\ref{fig:time}b),
a średnio są gorsze od PCA (0,524 i 0,567). Na dużym zbiorze najwięcej
kosztuje jednak detektor: OCSVM w danych oryginalnych Campaign uczy się
i ocenia średnio 95~s, IForest 0,85~s (aneks~\ref{app:full}).

\FloatBarrier
\subsection{Porównanie z zadaniem 1}

Tu zestawiono wyniki zadań 1 i 2 na wspólnym zbiorze danych i sprawdzono,
ile daje zmiana reprezentacji.

\begin{table}[H]
  \centering\small
  \caption{PR-AUC detektorów w każdej reprezentacji. Zad.~1 -- najlepsze
  ustawienie z zadania~1 (tylko Cardio jest wspólny). Zysk -- najlepsza
  reprezentacja minus dane oryginalne.}
  \label{tab:t1}
  \fittab{task1_comparison}
\end{table}

\textbf{Cardio: IForest i LOF mają w danych oryginalnych te same wyniki co
w zadaniu~1 (0,60 i 0,31)} (tabela~\ref{tab:t1}). OCSVM jest słabszy (0,60
wobec 0,76), bo w zadaniu~1 walidacja wybrała jądro sigmoidalne, a tu
siatka obejmuje tylko RBF. UMAP z OCSVM (0,84) przewyższa jednak najlepszy
wynik zadania~1 na tym zbiorze. Zmiana reprezentacji poprawiła wynik w 8 z 12
par zbiór--detektor, a o ponad 0,1 PR-AUC w sześciu.

\FloatBarrier
\section{Odpowiedzi na pytania analityczne}

W tej sekcji krótko odpowiedziano na pięć pytań postawionych w zadaniu, na
podstawie wyników z sekcji~5.

\begin{enumerate}[label=\alph*)]
  \item \textbf{Czy redukcja wymiarowości zawsze poprawia detekcję?} Nie.
  PCA poprawiło wynik w 8 z 12 par zbiór--detektor, autoenkoder w 5, UMAP
  w 2. Redukcja pomaga, gdy pozostałe wymiary to głównie szum (WDBC, $k=1$)
  albo cech jest więcej niż obserwacji (Arrhythmia). Szkodzi, gdy usuwa
  kierunki, w których anomalie się odchylają (Arrhythmia, $k=2$), albo
  wyrównuje gęstość (UMAP z LOF).
  \item \textbf{Czy anomalie pozostają oddzielone od danych normalnych?}
  W PCA z dobrym $k$ -- tak, na WDBC i Cardio nawet lepiej niż w danych
  oryginalnych (separacja do 0,71 i 0,46). UMAP zmniejsza separację na
  trzech zbiorach (WDBC: 0,53 $\rightarrow$ 0,09--0,34). Na Campaign
  separacja jest niska w danych oryginalnych, PCA i UMAP (0,07--0,28),
  a w autoenkoderze nieco wyższa (0,22--0,32), co jednak nie poprawia
  detekcji.
  \item \textbf{Która reprezentacja jest najlepsza dla którego detektora?}
  LOF -- PCA (największy zysk, bo najbardziej cierpi na szum wymiarów).
  OCSVM -- PCA, a gdy anomalie tworzą spójną grupę -- UMAP z dużym nn
  (Cardio). Isolation Forest -- dane oryginalne lub PCA z małym $k$; zmiana
  reprezentacji dała mu najwyżej +0,04.
  \item \textbf{Czy najlepsza wizualnie separacja oznacza najlepszą
  detekcję?} Nie. Najbardziej efektowne obrazy UMAP (Campaign, Arrhythmia)
  dają słabą detekcję, a liczbowa separacja zgadza się z jakością tylko na
  dwóch zbiorach na czterech.
  \item \textbf{Jaki jest kompromis jakości i kosztu?} Najlepszy ma PCA
  (milisekundy, najczęściej poprawa). UMAP i autoenkoder są tysiące razy
  droższe i rzadko lepsze. Na dużych danych więcej kosztuje detektor niż
  reprezentacja (OCSVM na Campaign: 95~s).
\end{enumerate}

\FloatBarrier
\section{Ograniczenia i dalsze badania}

W tej sekcji wymieniono, co ogranicza pewność wniosków i co warto zbadać
dalej.

\begin{itemize}
  \item WDBC ma w walidacji i teście po 2--3 anomalie, więc wyniki między
  podziałami różnią się bardzo (odchylenie 0,13--0,39), mimo 30 podziałów.
  Przewaga PCA na WDBC (0,92 wobec 0,71) jest jednak większa od tego rozrzutu
  i widać ją też na walidacji.
  \item Dla Campaign użyto 5 podziałów i węższej siatki UMAP i OCSVM.
  Rozrzut między podziałami jest tam mały (0,01--0,04), ale węższa siatka mogła
  pominąć lepszy wariant UMAP.
  \item Autoenkoder na Campaign nie skończył uczenia w limicie 400 epok.
  Zbadano tylko jeden typ sieci i uczenie na danych z anomaliami.
  \item Dalsze badania: odporne PCA, mniej wrażliwe na anomalie w części
  uczącej; autoenkoder uczony na danych oczyszczonych przez inny detektor;
  łączenie detektorów z różnych reprezentacji (zadanie~3).
\end{itemize}

\FloatBarrier
\section{Wnioski}

Najważniejsze wnioski z zadania, każdy poparty liczbami z sekcji~5.

\begin{enumerate}
  \item \textbf{Reprezentacja ma duży wpływ, w obie strony.} Ten sam detektor
  zmienia PR-AUC o ponad 0,2 zależnie od reprezentacji (OCSVM na WDBC: 0,65
  w autoenkoderze, 0,93 w PCA; LOF na Cardio: 0,24 w UMAP, 0,58 w PCA).
  \item \textbf{PCA jest najlepszą reprezentacją}: średnio PR-AUC 0,617 wobec
  0,544 w danych oryginalnych, przy koszcie kilku milisekund. Najbardziej
  zyskują LOF i OCSVM (WDBC: 0,71 $\rightarrow$ 0,92).
  \item \textbf{Liczby składowych nie wybierać według zachowanej
  wariancji}: na WDBC 33\,\% wariancji daje lepszy wynik niż 100\,\%,
  a na Arrhythmii optimum to 32--47\,\%.
  \item \textbf{UMAP zwykle szkodzi}, bo wyrównuje gęstość (pogorszenie w 10
  z 12 par), ale z OCSVM może oddzielić spójną grupę anomalii (Cardio: 0,84
  -- najlepszy wynik na tym zbiorze). Używać go głównie do rysunków, a do
  detekcji tylko z OCSVM i wyborem na walidacji.
  \item \textbf{Autoenkoder nie uzasadnia swojego kosztu} na tych danych
  (tysiące razy wolniejszy od PCA), poza Arrhythmią, gdzie cech jest więcej
  niż obserwacji.
  \item \textbf{Ładny obraz 2D nie oznacza dobrej detekcji}, a separacja
  klas zgadza się z jakością tylko na dwóch zbiorach z czterech.
  \item \textbf{Gdy anomalie nie odstają od danych (Campaign), żadna
  reprezentacja nie pomaga} (PR-AUC najwyżej 0,30). Żadna para
  reprezentacja--detektor nie jest najlepsza wszędzie -- to przesłanka do
  łączenia detektorów w zadaniu~3.
\end{enumerate}

\appendix
\FloatBarrier
\section{Wyniki wszystkich konfiguracji}
\label{app:full}

Tabela podaje wyniki najlepszych ustawień (według walidacji) każdej pary
reprezentacja--detektor: średnie z części testowych, $\pm$ odchylenie
standardowe dla PR-AUC. Skróty: oryg. -- dane oryginalne, AE -- autoenkoder,
rekonstr. -- błąd rekonstrukcji, Prec. i Rec. -- Precision i Recall; $t_r$ --
czas przygotowania reprezentacji, $t_d$ -- czas uczenia i oceny detektora.

{\footnotesize
\setlength{\tabcolsep}{2pt}
\tab{configs_full}
}

\end{document}
